% Options for packages loaded elsewhere
\PassOptionsToPackage{unicode}{hyperref}
\PassOptionsToPackage{hyphens}{url}
\documentclass[
]{article}
\usepackage{xcolor}
\usepackage[margin=1in]{geometry}
\usepackage{amsmath,amssymb}
\setcounter{secnumdepth}{-\maxdimen} % remove section numbering
\usepackage{iftex}
\ifPDFTeX
  \usepackage[T1]{fontenc}
  \usepackage[utf8]{inputenc}
  \usepackage{textcomp} % provide euro and other symbols
\else % if luatex or xetex
  \usepackage{unicode-math} % this also loads fontspec
  \defaultfontfeatures{Scale=MatchLowercase}
  \defaultfontfeatures[\rmfamily]{Ligatures=TeX,Scale=1}
\fi
\usepackage{lmodern}
\ifPDFTeX\else
  % xetex/luatex font selection
\fi
% Use upquote if available, for straight quotes in verbatim environments
\IfFileExists{upquote.sty}{\usepackage{upquote}}{}
\IfFileExists{microtype.sty}{% use microtype if available
  \usepackage[]{microtype}
  \UseMicrotypeSet[protrusion]{basicmath} % disable protrusion for tt fonts
}{}
\makeatletter
\@ifundefined{KOMAClassName}{% if non-KOMA class
  \IfFileExists{parskip.sty}{%
    \usepackage{parskip}
  }{% else
    \setlength{\parindent}{0pt}
    \setlength{\parskip}{6pt plus 2pt minus 1pt}}
}{% if KOMA class
  \KOMAoptions{parskip=half}}
\makeatother
\usepackage{color}
\usepackage{fancyvrb}
\newcommand{\VerbBar}{|}
\newcommand{\VERB}{\Verb[commandchars=\\\{\}]}
\DefineVerbatimEnvironment{Highlighting}{Verbatim}{commandchars=\\\{\}}
% Add ',fontsize=\small' for more characters per line
\usepackage{framed}
\definecolor{shadecolor}{RGB}{248,248,248}
\newenvironment{Shaded}{\begin{snugshade}}{\end{snugshade}}
\newcommand{\AlertTok}[1]{\textcolor[rgb]{0.94,0.16,0.16}{#1}}
\newcommand{\AnnotationTok}[1]{\textcolor[rgb]{0.56,0.35,0.01}{\textbf{\textit{#1}}}}
\newcommand{\AttributeTok}[1]{\textcolor[rgb]{0.13,0.29,0.53}{#1}}
\newcommand{\BaseNTok}[1]{\textcolor[rgb]{0.00,0.00,0.81}{#1}}
\newcommand{\BuiltInTok}[1]{#1}
\newcommand{\CharTok}[1]{\textcolor[rgb]{0.31,0.60,0.02}{#1}}
\newcommand{\CommentTok}[1]{\textcolor[rgb]{0.56,0.35,0.01}{\textit{#1}}}
\newcommand{\CommentVarTok}[1]{\textcolor[rgb]{0.56,0.35,0.01}{\textbf{\textit{#1}}}}
\newcommand{\ConstantTok}[1]{\textcolor[rgb]{0.56,0.35,0.01}{#1}}
\newcommand{\ControlFlowTok}[1]{\textcolor[rgb]{0.13,0.29,0.53}{\textbf{#1}}}
\newcommand{\DataTypeTok}[1]{\textcolor[rgb]{0.13,0.29,0.53}{#1}}
\newcommand{\DecValTok}[1]{\textcolor[rgb]{0.00,0.00,0.81}{#1}}
\newcommand{\DocumentationTok}[1]{\textcolor[rgb]{0.56,0.35,0.01}{\textbf{\textit{#1}}}}
\newcommand{\ErrorTok}[1]{\textcolor[rgb]{0.64,0.00,0.00}{\textbf{#1}}}
\newcommand{\ExtensionTok}[1]{#1}
\newcommand{\FloatTok}[1]{\textcolor[rgb]{0.00,0.00,0.81}{#1}}
\newcommand{\FunctionTok}[1]{\textcolor[rgb]{0.13,0.29,0.53}{\textbf{#1}}}
\newcommand{\ImportTok}[1]{#1}
\newcommand{\InformationTok}[1]{\textcolor[rgb]{0.56,0.35,0.01}{\textbf{\textit{#1}}}}
\newcommand{\KeywordTok}[1]{\textcolor[rgb]{0.13,0.29,0.53}{\textbf{#1}}}
\newcommand{\NormalTok}[1]{#1}
\newcommand{\OperatorTok}[1]{\textcolor[rgb]{0.81,0.36,0.00}{\textbf{#1}}}
\newcommand{\OtherTok}[1]{\textcolor[rgb]{0.56,0.35,0.01}{#1}}
\newcommand{\PreprocessorTok}[1]{\textcolor[rgb]{0.56,0.35,0.01}{\textit{#1}}}
\newcommand{\RegionMarkerTok}[1]{#1}
\newcommand{\SpecialCharTok}[1]{\textcolor[rgb]{0.81,0.36,0.00}{\textbf{#1}}}
\newcommand{\SpecialStringTok}[1]{\textcolor[rgb]{0.31,0.60,0.02}{#1}}
\newcommand{\StringTok}[1]{\textcolor[rgb]{0.31,0.60,0.02}{#1}}
\newcommand{\VariableTok}[1]{\textcolor[rgb]{0.00,0.00,0.00}{#1}}
\newcommand{\VerbatimStringTok}[1]{\textcolor[rgb]{0.31,0.60,0.02}{#1}}
\newcommand{\WarningTok}[1]{\textcolor[rgb]{0.56,0.35,0.01}{\textbf{\textit{#1}}}}
\usepackage{graphicx}
\makeatletter
\newsavebox\pandoc@box
\newcommand*\pandocbounded[1]{% scales image to fit in text height/width
  \sbox\pandoc@box{#1}%
  \Gscale@div\@tempa{\textheight}{\dimexpr\ht\pandoc@box+\dp\pandoc@box\relax}%
  \Gscale@div\@tempb{\linewidth}{\wd\pandoc@box}%
  \ifdim\@tempb\p@<\@tempa\p@\let\@tempa\@tempb\fi% select the smaller of both
  \ifdim\@tempa\p@<\p@\scalebox{\@tempa}{\usebox\pandoc@box}%
  \else\usebox{\pandoc@box}%
  \fi%
}
% Set default figure placement to htbp
\def\fps@figure{htbp}
\makeatother
\setlength{\emergencystretch}{3em} % prevent overfull lines
\providecommand{\tightlist}{%
  \setlength{\itemsep}{0pt}\setlength{\parskip}{0pt}}
\usepackage{bookmark}
\IfFileExists{xurl.sty}{\usepackage{xurl}}{} % add URL line breaks if available
\urlstyle{same}
\hypersetup{
  pdftitle={Template},
  pdfauthor={Studentnames and studentnumbers here},
  hidelinks,
  pdfcreator={LaTeX via pandoc}}

\title{Template}
\author{Studentnames and studentnumbers here}
\date{2026-06-17}

\begin{document}
\maketitle

\section{Set-up your environment}\label{set-up-your-environment}

\begin{Shaded}
\begin{Highlighting}[]
\FunctionTok{require}\NormalTok{(tidyverse)}
\end{Highlighting}
\end{Shaded}

\begin{verbatim}
## Loading required package: tidyverse
\end{verbatim}

\begin{verbatim}
## -- Attaching core tidyverse packages ------------------------ tidyverse 2.0.0 --
## v dplyr     1.2.1     v readr     2.2.0
## v forcats   1.0.1     v stringr   1.6.0
## v ggplot2   4.0.3     v tibble    3.3.1
## v lubridate 1.9.5     v tidyr     1.3.2
## v purrr     1.2.2     
## -- Conflicts ------------------------------------------ tidyverse_conflicts() --
## x dplyr::filter() masks stats::filter()
## x dplyr::lag()    masks stats::lag()
## i Use the conflicted package (<http://conflicted.r-lib.org/>) to force all conflicts to become errors
\end{verbatim}

\begin{Shaded}
\begin{Highlighting}[]
\FunctionTok{require}\NormalTok{(rmarkdown)}
\end{Highlighting}
\end{Shaded}

\begin{verbatim}
## Loading required package: rmarkdown
\end{verbatim}

\begin{Shaded}
\begin{Highlighting}[]
\FunctionTok{require}\NormalTok{(yaml)}
\end{Highlighting}
\end{Shaded}

\begin{verbatim}
## Loading required package: yaml
\end{verbatim}

\begin{Shaded}
\begin{Highlighting}[]
\FunctionTok{library}\NormalTok{(dplyr)}
\FunctionTok{library}\NormalTok{(systemfonts)}
\end{Highlighting}
\end{Shaded}

\section{Title Page}\label{title-page}

Changes in nitrogen emissions in the Netherlands

Group project by: Naut Leuverink, Lucas Wolffenbuttel, Mike Ottervanger,
Otto Ruyten, Nicolaas van Nispen en Robin Koldenhof

4.3

Fatima Amankour

\section{Part 1 - Identify a Social
Problem}\label{part-1---identify-a-social-problem}

Use APA referencing throughout your document.
\href{https://www.mendeley.com/guides/apa-citation-guide/}{Here's a link
to some explanation.}

\subsection{1.1 Describe the Social
Problem}\label{describe-the-social-problem}

Include the following:

\begin{itemize}
\item
  Why is this relevant?
\item
  Nitrogen emissions represent a critical public health crisis in the
  Netherlands. Elevated atmospheric concentrations lead to the
  production of particulate manner and ozone, which are closely
  associated with serious health problems such as worsening asthma,
  long-term cardiovascular problems, and diminished lung function.
\item
  This report measures the correlation between public health concerns
  and regional nitrogen levels over almost 20 years, from 2005 to 2023.
  By analyzing these temporal and spatial patterns, we want to pinpoint
  pollution hotspots and offer a data-driven framework for focused
  reduction initiatives, like industrial restrictions or localized
  agriculture changes, to successfully lessen the burden on public
  health.
\end{itemize}

\section{Part 2 - Data Sourcing}\label{part-2---data-sourcing}

\subsection{2.1 Load in the data}\label{load-in-the-data}

Preferably from a URL, but if not, make sure to download the data and
store it in a shared location that you can load the data in from. Do not
store the data in a folder you include in the Github repository!

\begin{Shaded}
\begin{Highlighting}[]
\NormalTok{PolutionInGeneral}\OtherTok{\textless{}{-}} \FunctionTok{read.csv}\NormalTok{(}\StringTok{"data/Polution in general.csv"}\NormalTok{)}
\NormalTok{BA\_BronnenVanStikstof }\OtherTok{\textless{}{-}} \FunctionTok{read.csv}\NormalTok{(}\StringTok{"data/Bronnen van stikstof.csv"}\NormalTok{)}
\NormalTok{AA\_BronnenVanFijnstof }\OtherTok{\textless{}{-}} \FunctionTok{read.csv}\NormalTok{(}\StringTok{"data/Bronnen van Fijnstof.csv"}\NormalTok{)}
\NormalTok{JaargemiddeldeFijnstof }\OtherTok{\textless{}{-}} \FunctionTok{read.csv}\NormalTok{(}\StringTok{"data/Jaargemiddelde PM10 per Locatie 2015{-}2025.csv"}\NormalTok{, }\AttributeTok{sep =} \StringTok{";"}\NormalTok{)}
\NormalTok{Meetlocaties }\OtherTok{\textless{}{-}} \FunctionTok{read.csv}\NormalTok{(}\StringTok{"data/luchtmeetnet\_meetlocaties.csv"}\NormalTok{)}


\FunctionTok{saveRDS}\NormalTok{(PolutionInGeneral, }\StringTok{"data/PolutionInGeneral.rds"}\NormalTok{)}
\NormalTok{PolutionInGeneral }\OtherTok{\textless{}{-}} \FunctionTok{readRDS}\NormalTok{(}\StringTok{"data/PolutionInGeneral.rds"}\NormalTok{)}
\FunctionTok{saveRDS}\NormalTok{(BA\_BronnenVanStikstof, }\StringTok{"data/BA\_BronnenVanStikstof.rds"}\NormalTok{)}
\NormalTok{BA\_BronnenVanStikstof }\OtherTok{\textless{}{-}} \FunctionTok{readRDS}\NormalTok{(}\StringTok{"data/BA\_BronnenVanStikstof.rds"}\NormalTok{)}
\FunctionTok{saveRDS}\NormalTok{(AA\_BronnenVanFijnstof, }\StringTok{"data/AA\_BronnenVanFijnstof.rds"}\NormalTok{)}
\NormalTok{AA\_BronnenVanFijnstof }\OtherTok{\textless{}{-}} \FunctionTok{readRDS}\NormalTok{(}\StringTok{"data/AA\_BronnenVanFijnstof.rds"}\NormalTok{)}
\FunctionTok{saveRDS}\NormalTok{(JaargemiddeldeFijnstof, }\StringTok{"data/JaargemiddeldeFijnstof.rds"}\NormalTok{)}
\NormalTok{JaargemiddeldeFijnstof }\OtherTok{\textless{}{-}} \FunctionTok{readRDS}\NormalTok{(}\StringTok{"data/JaargemiddeldeFijnstof.rds"}\NormalTok{)}
\FunctionTok{saveRDS}\NormalTok{(Meetlocaties, }\StringTok{"data/Meetlocaties.rds"}\NormalTok{)}
\NormalTok{Meetlocaties }\OtherTok{\textless{}{-}} \FunctionTok{readRDS}\NormalTok{(}\StringTok{"data/Meetlocaties.rds"}\NormalTok{)}
\end{Highlighting}
\end{Shaded}

\subsection{2.2 Provide a short summary of the
dataset(s)}\label{provide-a-short-summary-of-the-datasets}

\begin{Shaded}
\begin{Highlighting}[]
\CommentTok{\#head(PolutionInGeneral)}
\end{Highlighting}
\end{Shaded}

This dataset contains 11 variables, each representing a different year.
For each year, data is provided on the levels of nitrogen, NMVOCs
(non-methane volatile organic compounds), sulfur oxides, and PM2.5
particulate matter observed in the air. By analyzing this data, we can
assess the severity of air pollution and identify trends in air quality
over time.

\begin{Shaded}
\begin{Highlighting}[]
\CommentTok{\#head(JaargemiddeldeFijnstof)}
\end{Highlighting}
\end{Shaded}

These datasets show particulate matter (fine dust) levels measured at
various monitoring stations across the Netherlands over the years 2005
till 2025. This provides a clear insight into the periods and regions
where particulate matter causes the greatest impact and pollution.

\begin{Shaded}
\begin{Highlighting}[]
\CommentTok{\#head(BA\_BronnenVanStikstof)}
\end{Highlighting}
\end{Shaded}

This dataset shows the sources of nitrogen measured over a time period
from 2005 till 2024

\begin{Shaded}
\begin{Highlighting}[]
\CommentTok{\#head(AA\_BronnenVanFijnstof)}
\end{Highlighting}
\end{Shaded}

This dataset shows the sources of fine dust measured over a time period
from 2005 till 2024

\begin{Shaded}
\begin{Highlighting}[]
\NormalTok{inline\_code }\OtherTok{=} \ConstantTok{TRUE}
\end{Highlighting}
\end{Shaded}

These are things that are usually included in the metadata of the
dataset. For your project, you need to provide us with the information
from your metadata that we need to understand your dataset of choice.

\subsection{2.3 Describe the type of variables
included}\label{describe-the-type-of-variables-included}

Most of the datasets include a duration of time most often 2005 till
2024, and a variable that may be a specific measuring station or a
source of polution.

\section{Part 3 - Quantifying}\label{part-3---quantifying}

\subsection{3.1 Data cleaning}\label{data-cleaning}

Say we want to include only larger distances (above 2) in our dataset,
we can filter for this.

\begin{Shaded}
\begin{Highlighting}[]
\CommentTok{\#Laatste rij was .cbs, die hebben we verwijderd}
\NormalTok{AA\_BronnenVanFijnstof}\OtherTok{\textless{}{-}}\NormalTok{ AA\_BronnenVanFijnstof [}\SpecialCharTok{{-}}\FunctionTok{nrow}\NormalTok{(AA\_BronnenVanFijnstof),]}

\NormalTok{AA\_BronnenVanFijnstof }\OtherTok{\textless{}{-}}\NormalTok{ AA\_BronnenVanFijnstof }\SpecialCharTok{\%\textgreater{}\%}
  \FunctionTok{mutate}\NormalTok{(}\FunctionTok{across}\NormalTok{(}\FunctionTok{where}\NormalTok{(is.character),}\SpecialCharTok{\textasciitilde{}} \FunctionTok{as.numeric}\NormalTok{(}\FunctionTok{gsub}\NormalTok{(}\StringTok{","}\NormalTok{,}\StringTok{"."}\NormalTok{, .x))))}

\CommentTok{\#Nieuwe variabele gemaakt }
\NormalTok{AA\_BronnenVanFijnstof}\SpecialCharTok{$}\NormalTok{Totaal }\OtherTok{\textless{}{-}}\NormalTok{ AA\_BronnenVanFijnstof}\SpecialCharTok{$}\NormalTok{Landbouw..bosbouw.en.visserij..mln.kg.}\SpecialCharTok{+}\NormalTok{ AA\_BronnenVanFijnstof}\SpecialCharTok{$}\NormalTok{Particuliere.huishoudens..mln.kg.}\SpecialCharTok{+}\NormalTok{ AA\_BronnenVanFijnstof}\SpecialCharTok{$}\NormalTok{Energiesector..mln.kg.}\SpecialCharTok{+}\NormalTok{ AA\_BronnenVanFijnstof}\SpecialCharTok{$}\NormalTok{Industrie.en.afval..mln.kg.}\SpecialCharTok{+}\NormalTok{ AA\_BronnenVanFijnstof}\SpecialCharTok{$}\NormalTok{Vervoer..mln.kg.}\SpecialCharTok{+}\NormalTok{ AA\_BronnenVanFijnstof}\SpecialCharTok{$}\NormalTok{Diensten..water.en.bouw..mln.kg.}
\end{Highlighting}
\end{Shaded}

\begin{Shaded}
\begin{Highlighting}[]
\NormalTok{BA\_BronnenVanStikstof}\OtherTok{\textless{}{-}}\NormalTok{ BA\_BronnenVanStikstof [}\SpecialCharTok{{-}}\FunctionTok{nrow}\NormalTok{(BA\_BronnenVanStikstof),]}

\NormalTok{BA\_BronnenVanStikstof }\OtherTok{\textless{}{-}}\NormalTok{ BA\_BronnenVanStikstof }\SpecialCharTok{\%\textgreater{}\%}
  \FunctionTok{mutate}\NormalTok{(}\FunctionTok{across}\NormalTok{(}\FunctionTok{where}\NormalTok{(is.character),}\SpecialCharTok{\textasciitilde{}} \FunctionTok{as.numeric}\NormalTok{(}\FunctionTok{gsub}\NormalTok{(}\StringTok{","}\NormalTok{,}\StringTok{"."}\NormalTok{, .x))))}


\NormalTok{BA\_BronnenVanStikstof}\SpecialCharTok{$}\NormalTok{Totaal }\OtherTok{\textless{}{-}}\NormalTok{ BA\_BronnenVanStikstof}\SpecialCharTok{$}\NormalTok{Landbouw..bosbouw.en.visserij..mln.kg.}\SpecialCharTok{+}\NormalTok{ BA\_BronnenVanStikstof}\SpecialCharTok{$}\NormalTok{Particuliere.huishoudens..mln.kg.}\SpecialCharTok{+}\NormalTok{ BA\_BronnenVanStikstof}\SpecialCharTok{$}\NormalTok{Energiesector..mln.kg.}\SpecialCharTok{+}\NormalTok{ BA\_BronnenVanStikstof}\SpecialCharTok{$}\NormalTok{Industrie.en.afval..mln.kg.}\SpecialCharTok{+}\NormalTok{ BA\_BronnenVanStikstof}\SpecialCharTok{$}\NormalTok{Vervoer..mln.kg.}\SpecialCharTok{+}\NormalTok{ BA\_BronnenVanStikstof}\SpecialCharTok{$}\NormalTok{Diensten..water.en.bouw..mln.kg.}
\end{Highlighting}
\end{Shaded}

\subsubsection{Cleans the Meetlocaties dataset by splitting all data
into separate columns and removing the first four
rows.}\label{cleans-the-meetlocaties-dataset-by-splitting-all-data-into-separate-columns-and-removing-the-first-four-rows.}

\begin{Shaded}
\begin{Highlighting}[]
\CommentTok{\# We overschrijven \textquotesingle{}Meetlocaties\textquotesingle{} direct zodat de oude versie verdwijnt}
\NormalTok{Meetlocaties }\OtherTok{\textless{}{-}}\NormalTok{ Meetlocaties }\SpecialCharTok{\%\textgreater{}\%}
  \CommentTok{\# 1. Haal de rijen met een \textquotesingle{}\#\textquotesingle{} direct weg}
  \FunctionTok{filter}\NormalTok{(}\SpecialCharTok{!}\FunctionTok{str\_starts}\NormalTok{(.[[}\DecValTok{1}\NormalTok{]], }\StringTok{"\#"}\NormalTok{)) }\SpecialCharTok{\%\textgreater{}\%}
  
  \CommentTok{\# 2. Splits direct op de puntkomma\textquotesingle{}s naar losse kolommen}
  \FunctionTok{separate\_wider\_delim}\NormalTok{(}
    \AttributeTok{cols =} \DecValTok{1}\NormalTok{, }
    \AttributeTok{delim =} \StringTok{";"}\NormalTok{, }
    \AttributeTok{names =} \FunctionTok{c}\NormalTok{(}\StringTok{"meetlocatie\_id"}\NormalTok{, }\StringTok{"bron\_id"}\NormalTok{, }\StringTok{"meetlocatie\_naam"}\NormalTok{, }
              \StringTok{"meetlocatie\_plaatsnaam"}\NormalTok{, }\StringTok{"breedtegraad"}\NormalTok{, }\StringTok{"lengtegraad"}\NormalTok{, }
              \StringTok{"hoogte"}\NormalTok{, }\StringTok{"meetlocatie\_begindatumtijd"}\NormalTok{, }\StringTok{"meetlocatie\_einddatumtijd"}\NormalTok{)}
\NormalTok{  )}
\end{Highlighting}
\end{Shaded}

\subsubsection{Puts Meetlocaties \& JaargemiddeldeFijnstof together the
data gets a clear
location}\label{puts-meetlocaties-jaargemiddeldefijnstof-together-the-data-gets-a-clear-location}

\begin{Shaded}
\begin{Highlighting}[]
\FunctionTok{library}\NormalTok{(tidyverse)}

\CommentTok{\# 1. Zorg dat de koppelnaam in JaargemiddeldeFijnstof exact matcht}
\FunctionTok{names}\NormalTok{(JaargemiddeldeFijnstof)[}\DecValTok{1}\NormalTok{] }\OtherTok{\textless{}{-}} \StringTok{"meetlocatie\_id"}

\CommentTok{\# 2. Voer de koppeling uit en pak flexibel alle kolommen mee}
\NormalTok{Fijnstof\_Met\_Locaties }\OtherTok{\textless{}{-}}\NormalTok{ JaargemiddeldeFijnstof }\SpecialCharTok{\%\textgreater{}\%}
  \CommentTok{\# Koppel de schone Meetlocaties tabel eraan vast}
  \FunctionTok{left\_join}\NormalTok{(Meetlocaties, }\AttributeTok{by =} \StringTok{"meetlocatie\_id"}\NormalTok{) }\SpecialCharTok{\%\textgreater{}\%}
  
  \CommentTok{\# De nieuwe indeling: we zetten de handige kolommen vooraan, }
  \CommentTok{\# en met \textquotesingle{}everything()\textquotesingle{} komen ALLE overige kolommen (de jaren) er automatisch achteraan!}
  \FunctionTok{select}\NormalTok{(}
    \AttributeTok{Plaatsnaam =}\NormalTok{ meetlocatie\_plaatsnaam,}
    \AttributeTok{Breedtegraad =}\NormalTok{ breedtegraad,}
    \AttributeTok{Lengtegraad =}\NormalTok{ lengtegraad,}
    \FunctionTok{everything}\NormalTok{(),          }\CommentTok{\# Neemt automatisch alle overige kolommen/jaren mee}
    \SpecialCharTok{{-}}\NormalTok{meetlocatie\_id,       }\CommentTok{\# Haalt de oude NL{-}code netjes weg}
    \SpecialCharTok{{-}}\NormalTok{bron\_id,              }\CommentTok{\# Ruimt direct overbodige kolommen uit Meetlocaties op}
    \SpecialCharTok{{-}}\NormalTok{meetlocatie\_naam,}
    \SpecialCharTok{{-}}\NormalTok{hoogte,}
    \SpecialCharTok{{-}}\NormalTok{meetlocatie\_begindatumtijd,}
    \SpecialCharTok{{-}}\NormalTok{meetlocatie\_einddatumtijd}
\NormalTok{  )}
\end{Highlighting}
\end{Shaded}

\subsubsection{Makes Fijnstof\_Met\_Locaties
numeric}\label{makes-fijnstof_met_locaties-numeric}

\begin{Shaded}
\begin{Highlighting}[]
\NormalTok{Fijnstof\_Met\_Locaties }\OtherTok{\textless{}{-}}\NormalTok{ Fijnstof\_Met\_Locaties }\SpecialCharTok{\%\textgreater{}\%}
  \FunctionTok{mutate}\NormalTok{(}
    \CommentTok{\# 1. Coördinaten omzetten naar numeriek}
    \AttributeTok{Breedtegraad =} \FunctionTok{as.numeric}\NormalTok{(Breedtegraad),}
    \AttributeTok{Lengtegraad =} \FunctionTok{as.numeric}\NormalTok{(Lengtegraad)}
\NormalTok{  ) }\SpecialCharTok{\%\textgreater{}\%}
  \FunctionTok{mutate}\NormalTok{(}
    \CommentTok{\# 2. Alle jaarkolommen (die beginnen met X en daarna 4 cijfers) in één keer opschonen.}
    \CommentTok{\# Dit haalt de sterretjes, spaties en streepjes weg en maakt er getallen van.}
    \FunctionTok{across}\NormalTok{(}
      \AttributeTok{.cols =} \FunctionTok{matches}\NormalTok{(}\StringTok{"\^{}X}\SpecialCharTok{\textbackslash{}\textbackslash{}}\StringTok{d\{4\}$"}\NormalTok{), }
      \AttributeTok{.fns =} \SpecialCharTok{\textasciitilde{}} \FunctionTok{as.numeric}\NormalTok{(}\FunctionTok{str\_replace\_all}\NormalTok{(.x, }\StringTok{"([*]|}\SpecialCharTok{\textbackslash{}\textbackslash{}}\StringTok{s|{-})"}\NormalTok{, }\StringTok{""}\NormalTok{))}
\NormalTok{    )}
\NormalTok{  )}
\end{Highlighting}
\end{Shaded}

2 datasets naast elkaar gezet, en dat per jaar vergeleken, dataset
fijnstof gaat tot 2024 en dataset stikstof gaat tot 2023, daar moest ik
nog een aanpassing in maken.

Please use a separate `R block' of code for each type of cleaning. So,
e.g.~one for missing values, a new one for removing unnecessary
variables etc.

\subsection{3.2 Generate necessary
variables}\label{generate-necessary-variables}

\subsubsection{Creates total variables}\label{creates-total-variables}

\begin{Shaded}
\begin{Highlighting}[]
\CommentTok{\# De bestaande kolom \textquotesingle{}Totaal\textquotesingle{} wordt gebruikt.}
\CommentTok{\# De totalen van fijnstof en stikstof worden gecombineerd op basis van het jaar.}
\CommentTok{\# Omdat stikstof tot en met 2023 loopt, loopt de gecombineerde tabel ook tot en met 2023.}

\NormalTok{AB\_Fijnstof\_Totaal }\OtherTok{\textless{}{-}} \FunctionTok{data.frame}\NormalTok{(}
  \AttributeTok{Jaar =}\NormalTok{ AA\_BronnenVanFijnstof[, }\DecValTok{1}\NormalTok{],}
  \AttributeTok{Totaal\_fijnstof =}\NormalTok{ AA\_BronnenVanFijnstof}\SpecialCharTok{$}\NormalTok{Totaal}
\NormalTok{)}

\NormalTok{BB\_Stikstof\_Totaal }\OtherTok{\textless{}{-}} \FunctionTok{data.frame}\NormalTok{(}
  \AttributeTok{Jaar =}\NormalTok{ BA\_BronnenVanStikstof[, }\DecValTok{1}\NormalTok{],}
  \AttributeTok{Totaal\_stikstof =}\NormalTok{ BA\_BronnenVanStikstof}\SpecialCharTok{$}\NormalTok{Totaal}
\NormalTok{)}

\NormalTok{totale\_data }\OtherTok{\textless{}{-}} \FunctionTok{merge}\NormalTok{(}
\NormalTok{  AB\_Fijnstof\_Totaal,}
\NormalTok{  BB\_Stikstof\_Totaal,}
  \AttributeTok{by =} \StringTok{"Jaar"}
\NormalTok{)}
\end{Highlighting}
\end{Shaded}

\subsubsection{DIstancing data from the randstad from the
Non-randstad}\label{distancing-data-from-the-randstad-from-the-non-randstad}

\begin{Shaded}
\begin{Highlighting}[]
\CommentTok{\# Maak eerst de basisregio{-}dataset aan (als je dat nog niet had gedaan)}
\NormalTok{randstad\_plaatsen }\OtherTok{\textless{}{-}} \FunctionTok{c}\NormalTok{(}\StringTok{"Amsterdam"}\NormalTok{, }\StringTok{"Rotterdam"}\NormalTok{, }\StringTok{"Den Haag"}\NormalTok{, }\StringTok{"Utrecht"}\NormalTok{, }
                       \StringTok{"Haarlem"}\NormalTok{, }\StringTok{"Zaanstad"}\NormalTok{, }\StringTok{"Leiden"}\NormalTok{, }\StringTok{"Dordrecht"}\NormalTok{, }
                       \StringTok{"Almere"}\NormalTok{, }\StringTok{"Amersfoort"}\NormalTok{, }\StringTok{"Delft"}\NormalTok{, }\StringTok{"Schiedam"}\NormalTok{)}

\NormalTok{Fijnstof\_Regio }\OtherTok{\textless{}{-}}\NormalTok{ Fijnstof\_Met\_Locaties }\SpecialCharTok{\%\textgreater{}\%}
  \FunctionTok{mutate}\NormalTok{(}\AttributeTok{Regio =} \FunctionTok{if\_else}\NormalTok{(Plaatsnaam }\SpecialCharTok{\%in\%}\NormalTok{ randstad\_plaatsen, }\StringTok{"Randstad"}\NormalTok{, }\StringTok{"Other\_regions"}\NormalTok{))}

\CommentTok{\# {-}{-}{-} HIER SPLITSEN WE ZE IN TWEE LOSSE DATASETS {-}{-}{-}}
\NormalTok{Fijnstof\_Randstad      }\OtherTok{\textless{}{-}}\NormalTok{ Fijnstof\_Regio }\SpecialCharTok{\%\textgreater{}\%} \FunctionTok{filter}\NormalTok{(Regio }\SpecialCharTok{==} \StringTok{"Randstad"}\NormalTok{)}
\NormalTok{Fijnstof\_Niet\_Randstad }\OtherTok{\textless{}{-}}\NormalTok{ Fijnstof\_Regio }\SpecialCharTok{\%\textgreater{}\%} \FunctionTok{filter}\NormalTok{(Regio }\SpecialCharTok{==} \StringTok{"Other\_regions"}\NormalTok{)}
\end{Highlighting}
\end{Shaded}

\subsection{3.3 Create line graphs}\label{create-line-graphs}

\subsubsection{Simple graph about the general emission of particulate
matter
(fijnstof)}\label{simple-graph-about-the-general-emission-of-particulate-matter-fijnstof}

This graph shows the evolution of total particulate matter emissions
over the years from 2005 to 2025, where you can see that the total
particulate matter emissions have almost halved over the years.

\begin{Shaded}
\begin{Highlighting}[]
\FunctionTok{library}\NormalTok{(ggplot2)}

\FunctionTok{ggplot}\NormalTok{(AA\_BronnenVanFijnstof, }\FunctionTok{aes}\NormalTok{(}\AttributeTok{x =}\NormalTok{ X., }\AttributeTok{y =}\NormalTok{ Totaal)) }\SpecialCharTok{+}
  \FunctionTok{geom\_line}\NormalTok{(}\AttributeTok{linewidth =} \FloatTok{1.2}\NormalTok{) }\SpecialCharTok{+}
  \FunctionTok{geom\_point}\NormalTok{(}\AttributeTok{size =} \DecValTok{2}\NormalTok{) }\SpecialCharTok{+}
  \FunctionTok{labs}\NormalTok{(}
    \AttributeTok{title =} \StringTok{"Total particulate matter emissions over the years"}\NormalTok{,}
    \AttributeTok{x =} \StringTok{"Years"}\NormalTok{,}
    \AttributeTok{y =} \StringTok{"Particulate matter (mln kg)"}
\NormalTok{  ) }\SpecialCharTok{+}
  \FunctionTok{theme\_minimal}\NormalTok{()}
\end{Highlighting}
\end{Shaded}

\pandocbounded{\includegraphics[keepaspectratio]{Template_Assignment_files/Template_Assignment_files/figure-latex/unnamed-chunk-11-1.pdf}}

\subsubsection{Simple graph about the general emission of
nitrogen}\label{simple-graph-about-the-general-emission-of-nitrogen}

\begin{Shaded}
\begin{Highlighting}[]
\FunctionTok{library}\NormalTok{(ggplot2)}

\FunctionTok{ggplot}\NormalTok{(BA\_BronnenVanStikstof, }\FunctionTok{aes}\NormalTok{(}\AttributeTok{x =}\NormalTok{ X., }\AttributeTok{y =}\NormalTok{ Totaal)) }\SpecialCharTok{+}
  \FunctionTok{geom\_line}\NormalTok{(}\AttributeTok{linewidth =} \FloatTok{1.2}\NormalTok{) }\SpecialCharTok{+}
  \FunctionTok{geom\_point}\NormalTok{(}\AttributeTok{size =} \DecValTok{2}\NormalTok{) }\SpecialCharTok{+}
  \FunctionTok{labs}\NormalTok{(}
    \AttributeTok{title =} \StringTok{"Total nitrogen emissions over the years"}\NormalTok{,}
    \AttributeTok{x =} \StringTok{"Years"}\NormalTok{,}
    \AttributeTok{y =} \StringTok{"Nitrogen (mln kg)"}
\NormalTok{  ) }\SpecialCharTok{+}
  \FunctionTok{theme\_minimal}\NormalTok{()}
\end{Highlighting}
\end{Shaded}

\pandocbounded{\includegraphics[keepaspectratio]{Template_Assignment_files/Template_Assignment_files/figure-latex/unnamed-chunk-12-1.pdf}}

\subsubsection{Detailed and Multifaceted Map of Nitrogen Emission
Origins}\label{detailed-and-multifaceted-map-of-nitrogen-emission-origins}

\begin{Shaded}
\begin{Highlighting}[]
\CommentTok{\# 1. Laad de benodigde bibliotheken}
\FunctionTok{library}\NormalTok{(tidyverse)}

\CommentTok{\# 2. Bereid de dataset voor: hernoem de eerste kolom naar "X."}
\FunctionTok{names}\NormalTok{(BA\_BronnenVanStikstof)[}\DecValTok{1}\NormalTok{] }\OtherTok{\textless{}{-}} \StringTok{"X."}

\CommentTok{\# 3. Zet de tabel om naar het lange formaat dat jouw ggplot verwacht}
\NormalTok{BA\_BronnenVanStikstof\_long }\OtherTok{\textless{}{-}}\NormalTok{ BA\_BronnenVanStikstof }\SpecialCharTok{\%\textgreater{}\%}
  \FunctionTok{pivot\_longer}\NormalTok{(}
    \AttributeTok{cols =} \SpecialCharTok{{-}}\NormalTok{X., }
    \AttributeTok{names\_to =} \StringTok{"Bron"}\NormalTok{, }
    \AttributeTok{values\_to =} \StringTok{"Uitstoot"}
\NormalTok{  ) }\SpecialCharTok{\%\textgreater{}\%}
  \CommentTok{\# BELANGRIJK: Filter de "Totaal"{-}kolom eruit zodat deze de grafiek niet vertekent}
  \FunctionTok{filter}\NormalTok{(Bron }\SpecialCharTok{!=} \StringTok{"Totaal"}\NormalTok{)}


\FunctionTok{ggplot}\NormalTok{(BA\_BronnenVanStikstof\_long,}
       \FunctionTok{aes}\NormalTok{(}\AttributeTok{x =}\NormalTok{ X.,                                    }\CommentTok{\# \textless{}{-}{-}{-} Werkt nu met de jaren uit je CSV}
           \AttributeTok{y =}\NormalTok{ Uitstoot, }
           \AttributeTok{fill =}\NormalTok{ Bron)) }\SpecialCharTok{+}
  \FunctionTok{geom\_area}\NormalTok{(}\AttributeTok{color =} \StringTok{"white"}\NormalTok{, }\AttributeTok{linewidth =} \FloatTok{0.2}\NormalTok{) }\SpecialCharTok{+}       \CommentTok{\# \textless{}{-}{-}{-} Creëert de witte scheidingslijntjes}
  \FunctionTok{scale\_fill\_brewer}\NormalTok{(}
  \AttributeTok{palette =} \StringTok{"Set3"}\NormalTok{,}
  \AttributeTok{labels =} \FunctionTok{c}\NormalTok{(}
    \StringTok{"Services, water and construction"}\NormalTok{,}
    \StringTok{"Energy sector"}\NormalTok{,}
    \StringTok{"Industry and waste"}\NormalTok{,}
    \StringTok{"Agriculture, forestry and fisheries"}\NormalTok{,}
    \StringTok{"Households"}\NormalTok{,}
    \StringTok{"Transport"}
\NormalTok{  )}
\NormalTok{) }\SpecialCharTok{+}
  \FunctionTok{labs}\NormalTok{(                                               }\CommentTok{\# \textless{}{-}{-}{-} Namen op de assen en titels}
    \AttributeTok{title =} \StringTok{"Nitrogen emissions per source"}\NormalTok{,    }
    \AttributeTok{x =} \StringTok{"Years"}\NormalTok{,}
   \AttributeTok{y =} \StringTok{"Emissions (million kg)"}\NormalTok{,}
    \AttributeTok{fill =} \StringTok{"Sources"}
\NormalTok{  ) }\SpecialCharTok{+}
  \FunctionTok{theme\_bw}\NormalTok{() }\SpecialCharTok{+}                                        \CommentTok{\# \textless{}{-}{-}{-} Grijze rasterlijnen inschakelen}
  \FunctionTok{theme}\NormalTok{(}
    \AttributeTok{panel.grid.major =} \FunctionTok{element\_line}\NormalTok{(}\AttributeTok{color =} \StringTok{"grey80"}\NormalTok{),}
    \AttributeTok{panel.grid.minor =} \FunctionTok{element\_line}\NormalTok{(}\AttributeTok{color =} \StringTok{"grey90"}\NormalTok{),}
    \AttributeTok{plot.title =} \FunctionTok{element\_text}\NormalTok{(}\AttributeTok{hjust =} \FloatTok{0.5}\NormalTok{, }\AttributeTok{face =} \StringTok{"bold"}\NormalTok{)}
\NormalTok{  )}
\end{Highlighting}
\end{Shaded}

\pandocbounded{\includegraphics[keepaspectratio]{Template_Assignment_files/Template_Assignment_files/figure-latex/unnamed-chunk-13-1.pdf}}

\subsubsection{Detailed and Multifaceted Map of Particulate matter
Origins}\label{detailed-and-multifaceted-map-of-particulate-matter-origins}

\begin{Shaded}
\begin{Highlighting}[]
\CommentTok{\# 1. Laad de benodigde bibliotheken}
\FunctionTok{library}\NormalTok{(tidyverse)}

\CommentTok{\# 2. Bereid de dataset voor: hernoem de eerste kolom naar "X."}
\FunctionTok{names}\NormalTok{(AA\_BronnenVanFijnstof)[}\DecValTok{1}\NormalTok{] }\OtherTok{\textless{}{-}} \StringTok{"X."}

\CommentTok{\# 3. Zet de tabel om naar het lange formaat dat jouw ggplot verwacht}
\NormalTok{AA\_BronnenVanFijnstof\_long }\OtherTok{\textless{}{-}}\NormalTok{ AA\_BronnenVanFijnstof }\SpecialCharTok{\%\textgreater{}\%}
  \FunctionTok{pivot\_longer}\NormalTok{(}
    \AttributeTok{cols =} \SpecialCharTok{{-}}\NormalTok{X., }
    \AttributeTok{names\_to =} \StringTok{"Bron"}\NormalTok{, }
    \AttributeTok{values\_to =} \StringTok{"Uitstoot"}
\NormalTok{  ) }\SpecialCharTok{\%\textgreater{}\%}
  \CommentTok{\# BELANGRIJK: Filter de "Totaal"{-}kolom eruit zodat deze de grafiek niet vertekent}
  \FunctionTok{filter}\NormalTok{(Bron }\SpecialCharTok{!=} \StringTok{"Totaal"}\NormalTok{)}


\FunctionTok{ggplot}\NormalTok{(AA\_BronnenVanFijnstof\_long,}
       \FunctionTok{aes}\NormalTok{(}\AttributeTok{x =}\NormalTok{ X.,                                    }\CommentTok{\# \textless{}{-}{-}{-} Werkt nu met de jaren uit je CSV}
           \AttributeTok{y =}\NormalTok{ Uitstoot, }
           \AttributeTok{fill =}\NormalTok{ Bron)) }\SpecialCharTok{+}
  \FunctionTok{geom\_area}\NormalTok{(}\AttributeTok{color =} \StringTok{"white"}\NormalTok{, }\AttributeTok{linewidth =} \FloatTok{0.2}\NormalTok{) }\SpecialCharTok{+}       \CommentTok{\# \textless{}{-}{-}{-} Creëert de witte scheidingslijntjes}
  \FunctionTok{scale\_fill\_brewer}\NormalTok{(}
  \AttributeTok{palette =} \StringTok{"Set3"}\NormalTok{,}
  \AttributeTok{labels =} \FunctionTok{c}\NormalTok{(}
    \StringTok{"Services, water and construction"}\NormalTok{,}
    \StringTok{"Energy sector"}\NormalTok{,}
    \StringTok{"Industry and waste"}\NormalTok{,}
    \StringTok{"Agriculture, forestry and fisheries"}\NormalTok{,}
    \StringTok{"Households"}\NormalTok{,}
    \StringTok{"Transport"}
\NormalTok{  )}
\NormalTok{) }\SpecialCharTok{+}
  \FunctionTok{labs}\NormalTok{(                                               }\CommentTok{\# \textless{}{-}{-}{-} Namen op de assen en titels}
    \AttributeTok{title =} \StringTok{"Particulate matter emissions per source"}\NormalTok{,    }
    \AttributeTok{x =} \StringTok{"Years"}\NormalTok{,}
    \AttributeTok{y =} \StringTok{"Emissions (million kg)"}\NormalTok{,}
    \AttributeTok{fill =} \StringTok{"Sources"}
\NormalTok{  ) }\SpecialCharTok{+}
  \FunctionTok{theme\_bw}\NormalTok{() }\SpecialCharTok{+}                                        \CommentTok{\# \textless{}{-}{-}{-} Grijze rasterlijnen inschakelen}
  \FunctionTok{theme}\NormalTok{(}
    \AttributeTok{panel.grid.major =} \FunctionTok{element\_line}\NormalTok{(}\AttributeTok{color =} \StringTok{"grey80"}\NormalTok{),}
    \AttributeTok{panel.grid.minor =} \FunctionTok{element\_line}\NormalTok{(}\AttributeTok{color =} \StringTok{"grey90"}\NormalTok{),}
    \AttributeTok{plot.title =} \FunctionTok{element\_text}\NormalTok{(}\AttributeTok{hjust =} \FloatTok{0.5}\NormalTok{, }\AttributeTok{face =} \StringTok{"bold"}\NormalTok{)}
\NormalTok{  )}
\end{Highlighting}
\end{Shaded}

\pandocbounded{\includegraphics[keepaspectratio]{Template_Assignment_files/Template_Assignment_files/figure-latex/unnamed-chunk-14-1.pdf}}

\subsection{3.4 Creating a heatmap}\label{creating-a-heatmap}

\begin{verbatim}
## Linking to GEOS 3.14.1, GDAL 3.5.3, PROJ 9.5.1; sf_use_s2() is TRUE
\end{verbatim}

\pandocbounded{\includegraphics[keepaspectratio]{Template_Assignment_files/Template_Assignment_files/figure-latex/visualise_map-1.pdf}}

Here you provide a description of why the plot above is relevant to your
specific social problem.

\subsection{3.5 Visualize sub-population
variation}\label{visualize-sub-population-variation}

\subsubsection{A boxplot that combines the randstad and the
rest}\label{a-boxplot-that-combines-the-randstad-and-the-rest}

\begin{Shaded}
\begin{Highlighting}[]
\FunctionTok{library}\NormalTok{(tidyverse)}

\CommentTok{\# We gebruiken de gecombineerde dataset \textquotesingle{}Fijnstof\_Regio\textquotesingle{}}
\FunctionTok{ggplot}\NormalTok{(}\AttributeTok{data =}\NormalTok{ Fijnstof\_Regio, }\FunctionTok{aes}\NormalTok{(}\AttributeTok{x =}\NormalTok{ Regio, }\AttributeTok{y =} \FunctionTok{as.numeric}\NormalTok{(}\StringTok{\textasciigrave{}}\AttributeTok{X2025}\StringTok{\textasciigrave{}}\NormalTok{), }\AttributeTok{fill =}\NormalTok{ Regio)) }\SpecialCharTok{+}
  
\CommentTok{\# 1. Teken de boxplots naast elkaar (voeg na.rm = TRUE toe)}
\FunctionTok{geom\_boxplot}\NormalTok{(}\AttributeTok{width =} \FloatTok{0.5}\NormalTok{, }\AttributeTok{alpha =} \FloatTok{0.7}\NormalTok{,}
             \AttributeTok{outlier.color =} \StringTok{"red"}\NormalTok{, }\AttributeTok{outlier.shape =} \DecValTok{16}\NormalTok{, }\AttributeTok{outlier.size =} \DecValTok{2}\NormalTok{,}
             \AttributeTok{na.rm =} \ConstantTok{TRUE}\NormalTok{) }\SpecialCharTok{+}

\CommentTok{\# 2. Voeg de individuele meetstations als stipjes toe (voeg na.rm = TRUE toe)}
\FunctionTok{geom\_jitter}\NormalTok{(}\AttributeTok{width =} \FloatTok{0.15}\NormalTok{, }\AttributeTok{alpha =} \FloatTok{0.4}\NormalTok{, }\AttributeTok{color =} \StringTok{"grey30"}\NormalTok{, }
            \AttributeTok{na.rm =} \ConstantTok{TRUE}\NormalTok{) }\SpecialCharTok{+}
  
  \CommentTok{\# 3. Geef beide groepen een eigen herkenbare kleur}
  \FunctionTok{scale\_fill\_manual}\NormalTok{(}\AttributeTok{values =} \FunctionTok{c}\NormalTok{(}\StringTok{"Randstad"} \OtherTok{=} \StringTok{"coral"}\NormalTok{, }\StringTok{"Other\_regions"} \OtherTok{=} \StringTok{"skyblue"}\NormalTok{)) }\SpecialCharTok{+}
  
  \CommentTok{\# 4. Titels en labels aanpassen}
  \FunctionTok{labs}\NormalTok{(}
    \AttributeTok{title =} \StringTok{"Comparison Particulate matter emissions (2025)"}\NormalTok{,}
    \AttributeTok{subtitle =} \StringTok{"Randstad versus Rest of the country"}\NormalTok{,}
    \AttributeTok{y =} \StringTok{"PM2.5 Value (µg/m³)"}\NormalTok{,}
    \AttributeTok{x =} \StringTok{""} \CommentTok{\# De X{-}as laat nu netjes de categorienamen zien}
\NormalTok{  ) }\SpecialCharTok{+}
  
  \CommentTok{\# 5. De lay{-}out strak maken (de labels op de X{-}as laten we nu wél staan!)}
  \FunctionTok{theme\_minimal}\NormalTok{() }\SpecialCharTok{+}
  \FunctionTok{theme}\NormalTok{(}
    \AttributeTok{legend.position =} \StringTok{"none"}\NormalTok{, }\CommentTok{\# Legenda is overbodig omdat de X{-}as het al vertelt}
    \AttributeTok{plot.title =} \FunctionTok{element\_text}\NormalTok{(}\AttributeTok{face =} \StringTok{"bold"}\NormalTok{, }\AttributeTok{hjust =} \FloatTok{0.5}\NormalTok{),}
    \AttributeTok{plot.subtitle =} \FunctionTok{element\_text}\NormalTok{(}\AttributeTok{hjust =} \FloatTok{0.5}\NormalTok{)}
\NormalTok{  )}
\end{Highlighting}
\end{Shaded}

\pandocbounded{\includegraphics[keepaspectratio]{Template_Assignment_files/Template_Assignment_files/figure-latex/unnamed-chunk-15-1.pdf}}

By comparing the PM 2.5 particulate matter levels between the randstad
boxplot and the other regions boxplot, the boxplot comparison clearly
shows regional differences in Dutch air quality. This graphic directly
relates to the project's main social issue, which is examining the
effects of transportation, industrial infrastructure, and densely
populated areas on localized pollution and public health.

These plots demonstrate that randstad areas have a greater, more
concentrated load of particle matter, even though national data
indicates that industries like agriculture dominate nitrogen emissions
across rural regions. The randstad boxplot's larger median PM 2.5 value
highlights the unequal distribution of environmental effects and
presents air pollution as a pressing public health and urban development
issue.

\subsection{3.6 Event analysis}\label{event-analysis}

Analyze the relationship between two variables.

\begin{Shaded}
\begin{Highlighting}[]
\CommentTok{\# 1. Laad de benodigde bibliotheken}
\FunctionTok{library}\NormalTok{(tidyverse)}

\CommentTok{\# 2. Bereid de dataset voor: hernoem de eerste kolom naar "X."}
\FunctionTok{names}\NormalTok{(BA\_BronnenVanStikstof)[}\DecValTok{1}\NormalTok{] }\OtherTok{\textless{}{-}} \StringTok{"X."}

\CommentTok{\# (Optioneel: Zorg dat je jaartallen numeriek zijn, mocht de lijn niet werken)}
\CommentTok{\# BA\_BronnenVanStikstof$X. \textless{}{-} as.numeric(BA\_BronnenVanStikstof$X.)}

\CommentTok{\# 3. Zet de tabel om naar het lange formaat dat jouw ggplot verwacht}
\NormalTok{BA\_BronnenVanStikstof\_long }\OtherTok{\textless{}{-}}\NormalTok{ BA\_BronnenVanStikstof }\SpecialCharTok{\%\textgreater{}\%}
  \FunctionTok{pivot\_longer}\NormalTok{(}
    \AttributeTok{cols =} \SpecialCharTok{{-}}\NormalTok{X., }
    \AttributeTok{names\_to =} \StringTok{"Bron"}\NormalTok{, }
    \AttributeTok{values\_to =} \StringTok{"Uitstoot"}
\NormalTok{  ) }\SpecialCharTok{\%\textgreater{}\%}
  \CommentTok{\# BELANGRIJK: Filter de "Totaal"{-}kolom eruit zodat deze de grafiek niet vertekent}
  \FunctionTok{filter}\NormalTok{(Bron }\SpecialCharTok{!=} \StringTok{"Totaal"}\NormalTok{)}


\CommentTok{\# 4. Maak de plot}
\FunctionTok{ggplot}\NormalTok{(BA\_BronnenVanStikstof\_long,}
       \FunctionTok{aes}\NormalTok{(}\AttributeTok{x =}\NormalTok{ X.,                                    }\CommentTok{\# \textless{}{-}{-}{-} Werkt nu met de jaren uit je CSV}
           \AttributeTok{y =}\NormalTok{ Uitstoot, }
           \AttributeTok{fill =}\NormalTok{ Bron)) }\SpecialCharTok{+}
  \FunctionTok{geom\_area}\NormalTok{(}\AttributeTok{color =} \StringTok{"white"}\NormalTok{, }\AttributeTok{linewidth =} \FloatTok{0.2}\NormalTok{) }\SpecialCharTok{+}       \CommentTok{\# \textless{}{-}{-}{-} Creëert de witte scheidingslijntjes}
  
  \CommentTok{\# 5. lijn event analyse 2019 {-}{-}{-}}
  \FunctionTok{geom\_vline}\NormalTok{(}
  \AttributeTok{xintercept =} \DecValTok{2019}\NormalTok{,}
  \AttributeTok{color =} \StringTok{"darkred"}\NormalTok{,}
  \AttributeTok{linewidth =} \DecValTok{1}\NormalTok{,}
  \AttributeTok{linetype =} \StringTok{"dashed"}
\NormalTok{) }\SpecialCharTok{+}
\FunctionTok{annotate}\NormalTok{(}
  \StringTok{"text"}\NormalTok{,}
  \AttributeTok{x =} \DecValTok{2019}\NormalTok{,}
  \AttributeTok{y =} \DecValTok{285}\NormalTok{,}
  \AttributeTok{label =} \StringTok{"Nitrogen crisis}\SpecialCharTok{\textbackslash{}n}\StringTok{2019"}\NormalTok{,}
  \AttributeTok{color =} \StringTok{"darkred"}\NormalTok{,}
  \AttributeTok{angle =} \DecValTok{90}\NormalTok{,}
  \AttributeTok{vjust =} \SpecialCharTok{{-}}\FloatTok{0.5}
\NormalTok{) }\SpecialCharTok{+}
  \CommentTok{\# {-}{-}{-}{-}{-}{-}{-}{-}{-}{-}{-}{-}{-}{-}{-}{-}{-}{-}{-}{-}{-}{-}{-}{-}{-}{-}{-}{-}{-}{-}{-}{-}{-}{-}{-}{-}{-}{-}{-}{-}{-}{-}{-}{-}{-}{-}{-}{-}}
  
  \FunctionTok{scale\_fill\_brewer}\NormalTok{(}
  \AttributeTok{palette =} \StringTok{"Set3"}\NormalTok{,}
  \AttributeTok{labels =} \FunctionTok{c}\NormalTok{(}
    \StringTok{"Services, water and construction"}\NormalTok{,}
    \StringTok{"Energy sector"}\NormalTok{,}
    \StringTok{"Industry and waste"}\NormalTok{,}
    \StringTok{"Agriculture, forestry and fisheries"}\NormalTok{,}
    \StringTok{"Households"}\NormalTok{,}
    \StringTok{"Transport"}
\NormalTok{  )}
\NormalTok{) }\SpecialCharTok{+}
  \FunctionTok{labs}\NormalTok{(                                               }\CommentTok{\# \textless{}{-}{-}{-} Namen op de assen en titels}
    \AttributeTok{title =} \StringTok{"Nitrogen emissions per source"}\NormalTok{,  }
    \AttributeTok{x =} \StringTok{"Years"}\NormalTok{,}
    \AttributeTok{y =} \StringTok{"Emissions (million kg)"}\NormalTok{,}
    \AttributeTok{fill =} \StringTok{"Sources"}
\NormalTok{  ) }\SpecialCharTok{+}
  \FunctionTok{theme\_bw}\NormalTok{() }\SpecialCharTok{+}                                        \CommentTok{\# \textless{}{-}{-}{-} Grijze rasterlijnen inschakelen}
  \FunctionTok{theme}\NormalTok{(}
    \AttributeTok{panel.grid.major =} \FunctionTok{element\_line}\NormalTok{(}\AttributeTok{color =} \StringTok{"grey80"}\NormalTok{),}
    \AttributeTok{panel.grid.minor =} \FunctionTok{element\_line}\NormalTok{(}\AttributeTok{color =} \StringTok{"grey90"}\NormalTok{),}
    \AttributeTok{plot.title =} \FunctionTok{element\_text}\NormalTok{(}\AttributeTok{hjust =} \FloatTok{0.5}\NormalTok{, }\AttributeTok{face =} \StringTok{"bold"}\NormalTok{)}
\NormalTok{  )}
\end{Highlighting}
\end{Shaded}

\pandocbounded{\includegraphics[keepaspectratio]{Template_Assignment_files/Template_Assignment_files/figure-latex/analysis-1.pdf}}

The graph shows the development of nitrogen emissions from different
sources in the Netherlands between 2005 and 2024. Overall, emissions
decreased substantially over the observed period, with the largest
contribution consistently coming from the transport sector. Emissions
from industry, the energy sector, households, agriculture, and services
also show a downward trend.

The vertical dashed line indicates the Dutch nitrogen crisis in 2019.
After this point, emissions continued to decline across most sectors.
This trend coincides with the introduction of stricter environmental
measures and, shortly afterwards, the COVID-19 pandemic. However, based
on this graph alone, no direct causal relationship can be established
between these events and the observed decrease in emissions.

In the graph you see that there is a strong drop in nitrogen emissions
across multiple sources from 2019 onward. This can be explained by
multiple factors, such as new more strict emissions laws by the Dutch
government and also the EU. These laws include a lowering of the speed
limit to 100 km/h, causing a sharp drop in vehicle emissions. The laws
also include the freezing of a lot of infrastructure and other building
projects causing a drop in services and industry emissions. And then
from late 2019 and 2020 onward there was the COVID pandemic causing
lower emissions in almost all sectors.

\section{Part 4 - Discussion}\label{part-4---discussion}

\subsection{4.1 Discuss your findings}\label{discuss-your-findings}

\section{Part 5 - Reproducibility}\label{part-5---reproducibility}

\subsection{5.1 Github repository link}\label{github-repository-link}

Provide the link to your PUBLIC repository here:
\url{https://github.com/OLHHRuyten/Air-pollution-and-health-4.3-NL}

\subsection{5.2 Reference list}\label{reference-list}

Use APA referencing throughout your document. For the datasets:
AA\_BronnenVanFijnstof, BA\_BronnenVanStikstof and PolutionIngeneral are
from CBS:

Centraal Bureau voor de Statistiek. (2025, 2 september). Uitstoot
luchtverontreinigende stoffen lager dan norm voor 2030.
\url{https://www.cbs.nl/nl-nl/nieuws/2025/36/uitstoot-luchtverontreinigende-stoffen-lager-dan-norm-voor-2030}
Dataset JaargemiddeldeFijnstof Planbureau voor de Leefomgeving. (z.d.).
Atlas van de Regio. \url{https://themasites.pbl.nl/atlas-regio/}

\end{document}
